% Transcription — fonds Alexandre Grothendieck, shelfmark 100, batch 3 (pages 41-52).
% Established from the facsimile archives/batches/100/batch-03.pdf (a mirror of
% https://grothendieck.umontpellier.fr/100.pdf), per the skill
% transcribe-grothendieck. The batch file carries no cover sheet: PDF page N
% is page 40+N of the folder (the Montpellier footer prints « 41 » ... « 52 »),
% and the archivists' pencil numbers bottom-left agree wherever legible
% (« 41 », « 43 », « 45 », « 48 », « 49 », « 51 », « 52 »). The folder's last
% twelve pages; the batch is short because the folder is.
%
% Pass: Opus 5.5 (claude-opus-5-5), 2026-09-23 — first pass, unchecked against the pages by a human.
%
% The folder sits in the inventory group "69-89" « Géométrie et topologie
% combinatoire », dated 1976-[vers 1986], whose subject does not match this
% folder's; \dating{} carries the folder's own date verbatim, and nothing is
% concluded from the group.
%
% Eight of the twelve pages are transcribed. Absent:
%   - pages 41, 43 and 45, three leaves of a typed list « EGA IV 4e Errata
%     et Compléments » (45 carries that heading; 41 and 43 run on to
%     15.6.x, 8.10.3, 8.12.8, 12.1.1, 14.4.8, 14.4.10 and to a paragraph of
%     « autocritique » on Raynaud's henselisation and on « non ramifié » /
%     « net »). They carry pencil marks (brackets, « ? », a marginal
%     « 10.4.11.1 » on 43, diagonal cancellations on 43) which concern that
%     typescript only. Nothing in them bears on these notes; they are the
%     same kind of leaf batch 1 skipped at page 17. The typed text of 45 shows
%     through page 44 in mirror image, which is observed, not relied on.
%   - page 46, a blank leaf.
%
% What the batch is. Three runs, none dated.
%   42, 44  Blue ink, fast and heavily reworked. Page 42 opens on an NB that
%           refers back to the conditions (i)-(iv) closing batch 2 (page 40:
%           algebraisability and projectivity of \hat X over \hat S'): if
%           X_0 is an abelian scheme, the projectivity condition can be
%           dropped. Then « Remarques »: S' -> S a morphism of effective
%           descent, S loc. noeth.; X' flat and proper over S' with a
%           descent datum relative to S — is the datum effective? a) reduce
%           to the local rings O_{S,s}; b) to the completions; c) fragments.
%           Page 44 opens on « d) » (so continues a)-c) of page 42 as far as
%           the lettering shows) — the case where some S'' -> S' is
%           surjective with X projective over S', Pic^0 of the X'_n/S'_n
%           flat, the fibres without infinitesimal automorphisms: descent
%           after a finite étale extension T of Spec \hat O_{S,s}; then a
%           second « Remarque »-like paragraph with a polarised pair
%           (X', L') and a descent datum, (X'_{s'}, L'_{s'}) without
%           infinitesimal automorphisms, « Alors la donnée … est effective.
%           Argument par modules formels… ». Most connective prose on both
%           pages is lost and left \ill{}; the lower half of 42 is given
%           only in fragments, with a \note saying so.
%   47-50   A typescript, « Modèles de Néron et groupes p-divisibles »,
%           paginated (unnumbered first leaf), « - 2 - », « - 3 - »,
%           « - 4 - »; not signed, not dated, author not named on the leaves
%           (it speaks in the first person: « je sais maintenant prouver le
%           théorème de Tate », « séminaire de l'an dernier », « note sur les
%           modèles de Néron prop. 2.5 »). S a trait, generic point η:
%           Theorem — for G a Néron-type S-group scheme and H a p-divisible
%           group over S, every u_η : H_η -> G_η extends uniquely; Lemma 1
%           (reduction to the connected part M~ of a finite flat M), Lemma 2
%           (M x_S G / P smooth), proof of the theorem in the semi-stable
%           case (Tate twice) and in general (Weil restriction Π_{S'/S},
%           schematic closure G_1, « lissifié », Lemma 3: p^n on G_1
%           factors through G), a remark on the use of Tate's theorem, and
%           a Corollary (good reduction iff the p-divisible group extends).
%           Transcribed as typed text. Annotations: blue ink (the η the
%           typewriter could not strike, wavy underlines marking p-divisible
%           groups, the script 𝒢 and the Π, small corrections « dé »,
%           « un », « étale ») and pencil (marginal questions and
%           comments); both are given in the apparatus, the pencil ones as
%           \marginal{}. The folder title calls these « copies de tapuscrit
%           annoté »; whose hand the annotations are is not settled here.
%   51-52   Blue ink, fair. A « Lemme »: K a field, G an algebraic group, M a
%           p-divisible group, φ : M -> G; φ factors through a subgroup H
%           that is an extension of an abelian variety by a torus (under the
%           boxed hypothesis car K = p or K perfect, and G without subgroup
%           ≅ G_a); proof via the centralisers Z_n = Centr φ(n), reduction
%           to G commutative connected, then p^ν G for large ν. Page 52 draws
%           the consequence for a Néron model G over a trait: φ_K extends to
%           φ : M -> G, via the Néron model H of H_K and « le th. de Raynaud
%           pour H ». « On gagne ! » closes the folder.
%
% Notation fixed for the batch. The typescript marks p-divisible groups by a
% wavy underline (blue ink): set as \underset{\sim}{G}, \underset{\sim}{H};
% G_η etc. with the η he wrote in by hand; the blue curly G of the Néron model
% over S' as \mathcal{G}; M~, u~, v~ (typed tildes) as \widetilde; the Weil
% restriction as \prod_{S'/S}. On 51-52: Centr, VA, Z_n, p^ν G as written.
%
% Cross-references. Page 42's « (i) à (iv) » are the conditions at the end of
% batch 2 (page 40). No explicit page references otherwise. The typescript's
% Theorem is the statement page 52 proves by hand for G a Néron model; the
% folder title's « bonne réduction … via bonne réduction de Barsotti-Tate »
% is the typescript's Corollary (page 49-50).
\documentclass[11pt,a4paper]{article}
% Le préambule commun à toutes les transcriptions du fonds Grothendieck.
%
% Il tient en une idée : **l'incertitude se compose, elle ne se résout pas.**
% Une transcription qui lisse un mot douteux a détruit la seule chose pour
% laquelle on la faisait — un lecteur doit pouvoir savoir, à chaque endroit,
% ce qui était sur le feuillet et ce qui a été ajouté. Les six macros qui
% suivent sont donc l'appareil critique, et non de la décoration.
%
% Les mêmes commandes sont reconnues par `scripts/render.mjs`, qui produit la
% vue de lecture du site. Une macro ajoutée ici doit l'être aussi là-bas,
% faute de quoi le rendu échouera bruyamment — ce qui est voulu : un rendu
% silencieusement faux est pire qu'un rendu absent.

\ProvidesPackage{grothendieck}[2026/08/08 Transcriptions du fonds Grothendieck]

\usepackage{amsmath,amssymb,amsthm}
\usepackage{mathrsfs}
\usepackage{stmaryrd}
% Les diagrammes commutatifs sont partout dans ce fonds : c'est la langue
% même de la géométrie algébrique de Grothendieck, et une transcription qui
% les rendrait en prose ne transcrirait rien.
\usepackage{tikz-cd}
% Français par défaut — c'est la langue des feuillets — mais l'anglais est
% chargé aussi : la traduction et le résumé sont des documents anglais, et la
% typographie française (espaces avant les deux-points) y serait une faute.
% Ces documents basculent par \selectlanguage{english} en tête de corps.
\usepackage[english,french]{babel}
\usepackage{fontspec}
\usepackage{geometry}
\geometry{a4paper,margin=25mm,marginparwidth=28mm}
\usepackage{marginnote}
\usepackage{xcolor}
% ulem, not soul. `soul`'s \st and \ul are built on a character-by-character
% scan that XeLaTeX's font machinery breaks: the document fails with an
% undefined control sequence rather than degrading. ulem does the same two jobs
% and is Unicode-engine safe. `normalem` stops it hijacking \emph, which the
% transcriptions use for Grothendieck's own emphasis.
\usepackage[normalem]{ulem}
\usepackage{hyperref}
\hypersetup{colorlinks=true,linkcolor=black,urlcolor=black,pdfborder={0 0 0}}

% --- Métadonnées du lot -----------------------------------------------------
% Reprises telles quelles par le script de rendu, qui les affiche en tête de
% la vue de lecture. Elles fixent ce que le fichier prétend couvrir : sans
% elles, une transcription flotte sans point d'attache dans un fonds de seize
% mille feuillets.

\newcommand{\folder}[1]{\gdef\grothFolder{#1}}
\newcommand{\batch}[1]{\gdef\grothBatch{#1}}
\newcommand{\leaves}[2]{\gdef\grothFirst{#1}\gdef\grothLast{#2}}
\newcommand{\foldertitle}[1]{\gdef\grothFoldertitle{#1}}
\gdef\grothFolder{?}\gdef\grothBatch{?}\gdef\grothFirst{?}\gdef\grothLast{?}\gdef\grothFoldertitle{}

% --- Le résumé --------------------------------------------------------------
%
% Il ouvre la lecture modernisée, sous le titre « Résumé », et il est composé
% en petit. Ce n'est pas un ornement : c'est le seul passage du document qui
% ne soit pas des mathématiques, et un lecteur doit pouvoir le reconnaître
% avant de l'avoir lu — puis le sauter s'il connaît déjà le sujet, ou s'y
% arrêter s'il ne le connaît pas. Le filet qui le suit marque la reprise du
% corps.

\newenvironment{resume}
  {\small\setlength{\parskip}{0.45em}}
  {\par\medskip\hrule\medskip}

% --- Mots-clés ---------------------------------------------------------------
%
% En anglais, en clôture du résumé de la lecture modernisée : le vocabulaire
% moderne sous lequel chercher ce que les feuillets construisent. Le manifeste
% du site (`npm run manifest`) les extrait pour étiqueter la cote entière —
% cette ligne est la source unique des tags, il n'y a pas de fichier à côté.
\newcommand{\keywords}[1]{\par\smallskip\noindent
  {\footnotesize\textsc{Keywords} --- \textit{#1}}\par}

% --- L'appareil critique ----------------------------------------------------

% Le numéro de feuillet, en marge. C'est l'ancre qui permet de régler un
% désaccord sur une formule en regardant *un* feuillet plutôt qu'une chemise
% de six cents. Le site s'en sert aussi pour tourner les pages du fac-similé
% quand on fait défiler la transcription.
\newcommand{\leaf}[1]{%
  \marginnote{\footnotesize\color{gray!70}\textsf{#1}}%
  \label{leaf:#1}\ignorespaces}

% Illisible : on ne devine pas. Un mot inventé qui se lit comme les autres est
% le pire résultat possible.
\newcommand{\ill}{\textcolor{red!60!black}{[\,\ldots\,]}}

% Lecture douteuse : le mot est donné, et signalé comme incertain.
\newcommand{\uncertain}[1]{\uline{#1}}

% Ajout de l'éditeur — une lettre manquante, un mot sous-entendu.
\newcommand{\add}[1]{\textcolor{blue!50!black}{[#1]}}

% Biffé par Grothendieck lui-même. Ce qu'il a rayé fait partie du document :
% une rature dit souvent où la pensée a changé de direction.
\newcommand{\struck}[1]{\sout{#1}}

% Note du transcripteur, sur l'état matériel du feuillet.
\newcommand{\note}[1]{\par\smallskip{\small\color{gray!60!black}\textit{[#1]}}\par\smallskip}

% Note marginale de Grothendieck, distincte du corps du texte.
\newcommand{\marginal}[1]{\par\smallskip\noindent\hspace{1em}%
  {\small\color{gray!60!black}#1}\par\smallskip}

% --- Titre ------------------------------------------------------------------

\AtBeginDocument{%
  \begin{center}
    {\large\bfseries Fonds Alexandre Grothendieck --- cote n\textsuperscript{o}~\grothFolder}\\[2pt]
    {\itshape\grothFoldertitle}\\[4pt]
    {\small feuillets \grothFirst--\grothLast\ (lot \grothBatch)}\\[2pt]
    {\footnotesize\color{gray!60!black}Transcription établie d'après le fac-similé numérisé
     par l'Université de Montpellier.\\
     Les lectures incertaines sont soulignées, les illisibles notées [\,\ldots\,],
     les ajouts entre crochets.}
  \end{center}
  \vspace{1em}}

\endinput


\folder{100}
\batch{3}
\pages{41}{52}
\dating{[à partir de 1967-1968]}
\watermark{Édition de démonstration}
\foldertitle{Bonne réduction des variétés abéliennes via bonne réduction de Barsotti-Tate Tp (A) (base quelconque) : notes manuscrites (s.d.), copies de tapuscrit annoté (s.d.).}

\begin{document}

\section{[Descente des schémas propres et plats : remarques]}
\note{titre de l'éditeur, entre crochets : les feuillets 42 et 44 n'en portent
aucun. Le feuillet 42 enchaîne sur les conditions (i) à (iv) qui terminent le
lot 2 (page 40).}

\page{42}
\emph{N.B.} Si $X_0$ est un schéma abélien \struck{et} (ou seulement
\uncertain{pro-abélien} ?), alors ces conditions \struck{impliquent}
[\uncertain{équivalent à celle}] que $\hat{X}$ soit \uncertain{algébrisable},
et dans (i) à (iv) [\ill{} \ill{}] on peut supprimer la condition de
projectivité.
\marginal{?}
\note{« équivalent à celle » est écrit au-dessus du mot biffé ; un second
ajout interlinéaire, après « (i) à (iv) », ne se lit pas. Le « ? » est en
marge gauche, en face de la première ligne.}

\textbf{Remarques.} Soit $S' \to S$ un morphisme [\ill{} \ill{}]
\uncertain{de descente} effectif ($S$ loc.\ noeth.). \ill{} donc un morphisme
de descente effective pour les schémas relatifs plats.
\note{un ajout encadré de deux lignes surmonte « un morphisme » ; on n'y lit
que le mot « morphisme ».}

Soit $X'$ plat et propre \struck{\ill{}} sur $S'$, avec donnée de descente
relativement à $S$, on \struck{\ill{}} se demande si la donnée de descente est
effective.

a) \struck{\uncertain{Condition locale sur $S$}} \ill{} on peut se réduire aux
localisés $\mathcal{O}_{S,s}$ de $S$.

b) Moyennant une question de descente \struck{\ill{}} \uncertain{fid.} \ill{},
\uncertain{complète}, on est ramené au problème sur
$\hat{\mathcal{O}}_{S,s}$. \ill{} \ill{} \ill{} de $X_0$ \ill{} \ill{}
$\hat{\mathcal{O}}_{S,s}$ \ill{} \ill{} résiduel $k$ \ill{} \ill{} \ill{}
\uncertain{descendre} \ill{} les fibres \ill{} \ill{}

c) \ill{} \ill{} que les fibres de $X'/S'$ soient \ill{} \ill{}
\uncertain{infinitésimaux} \ill{} \struck{\uncertain{soit lisse sur $S'$}}
\ill{} \ill{} \ill{} \ill{} un \ill{} \ill{} \uncertain{complété} \ill{}.
Alors \ill{} un schéma formel $\hat{X}$ sur $\hat{S}$ qui descend $\hat{X}'$.
\marginal{\uncertain{Conditions pour} \ill{} \ill{} \uncertain{constantes}}
\note{la moitié inférieure du feuillet, à partir de b), est un premier jet
très rapide, surchargé de ratures et d'ajouts interlinéaires ; on n'en donne
que les fragments lisibles. La note marginale est à gauche, en face de c).}

\page{44}
d) Si on suppose [\uncertain{au lieu de la} \ill{}] qu'il existe
$S'' \to S'$ \struck{\uncertain{quasi-fini}} surjectif avec
$\bar{X}_{S'}$ \struck{\ill{}} projectif sur $S'$, alors on démontre que
\ill{} \ill{} pour tout $s \in S$, \struck{\ill{}} tels que pour
\ill{} $s' \in S'_s$,
\ill{} $\mathrm{Pic}^{0}_{X'_n/S'_n}$
($S'_n = \mathrm{Spec}\, \mathcal{O}_{S',s'}/\mathfrak{m}_{s'}^{n+1}$) \ill{}
plats, et \uncertain{que} les $X'_{s'}$ ($s' \in S'_s$) soient sans
automorphismes infinitésimaux, il existe une extension finie étale $T$ de
$\mathrm{Spec}\, \hat{\mathcal{O}}_{S,s}$ telle que la donnée de descente
\struck{\ill{}} sur $T'/T$ sur $X'_{T'}$ \ill{} \ill{} \ill{} \ill{} $X_T$
sur $T$ ; \ill{} \ill{}, \ill{} \ill{} \ill{} \ill{} que les
\struck{\ill{}} $X_{T_i}$, où les $T_i$ sont les \uncertain{comp.\
irréductibles} de $T$, \ill{} \ill{} $(T_i)$ \ill{}
\uncertain{projectivité} \ill{}).
\note{l'ajout « au lieu de la … » est écrit au-dessus de la ligne, puis
entouré ; le mot biffé après $S'' \to S'$ pourrait être « quasi-fini ». On
lit ensuite, écrit en bas de ligne, « (ce $T'$ …) », puis « $\mathrm{Spec}\,
\hat{\mathcal{O}}_{S,s}$ » avec son chapeau.}

\textbf{\uncertain{Remarques}.} Soit donc $S' \to S$ un morphisme
\uncertain{propre} \ill{} \uncertain{effectif}, avec $S$ loc.\ noeth.
Supposons qu'on ait sur $S'$ un schéma \uncertain{propre et} plat $X'/S'$, un
faisceau inversible $\underline{L}'$ sur $X'$, et une donnée de descente sur
$(X', \underline{L}')$ rel.\ à $S' \to S$.
\note{« propre et » est ajouté au-dessus de « plat ».}

\struck{Alors} Supposons de plus que pour tout $s' \in S'$,
$(X'_{s'}, \underline{L}'_{s'})$ n'ait pas d'automorphisme infinitésimal
\ill{} \ill{} \ill{} donnée de descente \uncertain{considérée} \ill{}.
Alors la donnée \ill{} \ill{} est effective. Argument par modules formels…

\section{Modèles de Néron et groupes $p$-divisibles}
\note{titre dactylographié et souligné en tête du feuillet 47, premier d'un
tapuscrit de quatre feuillets (47, puis « - 2 - » à « - 4 - » sur 48 à 50),
sans signature ni date. Le texte tapé est transcrit tel quel ; les
annotations à l'encre bleue (les $\eta$ que la machine ne portait pas, le
trait ondulé qui marque les groupes $p$-divisibles, rendu
$\underset{\sim}{G}$, la lettre $\mathcal{G}$ et le signe $\prod$) sont
intégrées au texte, les autres signalées ; les notes au crayon sont données
en « marginal ».}

\page{47}
Dans ce qui suit $S$ est un trait de point générique $\eta$ et de point fermé
$s$. Un $S$-groupe $p$-divisible $(G_n,\ n \in \mathbb{N})$ sera aussi
désigné par le symbole $\underset{\sim}{G}$. On \supplied{dé}signe par $p$ un
nombre premier fixé.
\note{la machine a tapé « séigne » ; « dé » est écrit à l'encre bleue au-dessus.}

\textbf{Théorème} : Soient $G$ un $S$-schéma en groupes, $S$-néronien, [dont
la fibre générique est extension d'une variété abélienne par un tore] et
$\underset{\sim}{H}$ un $S$-groupe $p$-divisible. Alors tout morphisme
$u_\eta : (\underset{\sim}{H})_\eta \to G_\eta$ se prolonge de manière unique
en un $S$-morphisme $u : \underset{\sim}{H} \to G$.
\marginal{[\,] \uncertain{inutile}}
\note{les crochets droits autour de « dont la fibre …\ tore » sont au
crayon, repris en marge sous la forme « [\,] » suivie d'un mot que l'on lit
« inutile » sans certitude.}

\textbf{Lemme 1} : Supposons $S$ hensélien. Soient $G$ un $S$-schéma en
groupes, $S$-néronien \struck{et}, $M$ un $S$-schéma en groupes fini et plat,
$\widetilde{M}$ sa partie connexe, \struck{$u_K$} $u_\eta : M_\eta \to G_\eta$
un homomorphisme. Alors, pour que $u_\eta$ se prolonge en un $S$-morphisme
$u : M \to G$, il faut et il suffit que $u_\eta \mid \widetilde{M}_\eta$ se
prolonge en un $S$-morphisme $\widetilde{u} : \widetilde{M} \to G$.

La nécessité est claire. Prouvons la suffisance. Considérons le groupe produit
$M \times_S G$ et le morphisme
\[
  v_\eta : M_\eta \times G_\eta \longrightarrow G_\eta, \qquad
  (m, g) \longmapsto u_\eta(m) + g .
\]
\marginal{$G$ commutatif ?}
Il nous suffit de montrer que $v_\eta$ se prolonge en un morphisme
$v : M \times_S G \to G$, car la restriction de $v$ au premier facteur $M$
fournira le prolongement de $u$. Soit $P_\eta$ le noyau de $v_\eta$ et $P$
l'adhérence schématique de $P_\eta$ dans $M \times_S G$. Le lemme 1 résulte de
la propriété néronienne de $G$ et du lemme suivant :
\note{« $G$ commutatif ? » est au crayon dans la marge gauche, en face de
« produit $M \times_S G$ ».}

\textbf{Lemme 2} : Le $S$-schéma en groupes $M \times_S G/P$ [existe et] est
lisse sur $S$.
\note{« existe et » est ajouté au crayon au-dessus de la ligne précédente,
relié par un trait à la place où il s'insère.}

Vu l'hypothèse faite sur $u_\eta$, le morphisme
$\widetilde{v}_\eta = v_\eta \mid \widetilde{M}_\eta \times G_\eta$ se
prolonge en un $S$-morphisme
\[
  \widetilde{v} : \widetilde{M} \times_S G \longrightarrow G, \qquad
  (\widetilde{m}, g) \longmapsto \widetilde{u}(\widetilde{m}) + g ,
\]
de sorte que l'on a la suite exacte :
\[
  (1) \qquad 0 \to \operatorname{Ker} \widetilde{v} \to \widetilde{M} \times_S G
  \to G \to 0 .
\]
\note{un trait de crayon en marge droite, en face de la suite (1).}

\page{48}
\marginal{\uncertain{en fait il y a} \ill{} : $\widetilde{M}$ !}
\note{note au crayon en haut du feuillet, à droite de « - 2 - », reliée par un
trait à la première ligne ; lecture très incertaine.}

Il en résulte que $\operatorname{Ker}(\widetilde{v})$ est plat et que
$\widetilde{M} \times_S G/\operatorname{Ker}(\widetilde{v})$ est lisse. Par
ailleurs, vu le choix de $\widetilde{M}$, Le groupe
$M \times_S G / \widetilde{M} \times_S G = M/\widetilde{M}$ est étale, donc
$M \times_S G/\operatorname{Ker}(\widetilde{v})$ est lisse, comme extension de
groupes lisses. On a évidemment $\operatorname{Ker}(\widetilde{v}) \subset P$,
de sorte que $M \times_S G/P$ est \struck{bien} lisse, comme quotient d'un
groupe lisse.

Démonstration du théorème. L'unicité du prolongement est claire. Comme la
propriété néronienne de $G$ est conservée par passage à l'hensélisé, on peut
supposer $S$ hensélien. Quitte alors à remplacer $\underset{\sim}{H}$ par sa
partie connexe, on peut, compte tenu du lemme 1, supposer
\struck{\ill{}} $\underset{\sim}{H}$ connexe.

a) Cas où $G$ est semi-stable. Soit $\underset{\sim}{G}_\eta$ le groupe
$p$-divisible défini par $G_\eta$ et soit $\underset{\sim}{G}'$ le
$S$-groupe $p$-divisible obtenu à partir du schéma en groupes $G^{\circ}$. On
sait que
$\underset{\sim}{G}''_\eta = \underset{\sim}{G}_\eta/\underset{\sim}{G}'_\eta$
est \add{un} groupe $p$-divisible \struck{non ramifié} [étale] sur $S$,
\struck{donc} [et] se prolonge en un $S$-groupe $p$-divisible,
$\underset{\sim}{G}''$, étale sur $S$. D'après Tate, le morphisme composé
$v_\eta : \underset{\sim}{H}_\eta \xrightarrow{u_\eta}
\underset{\sim}{G}_\eta \to \underset{\sim}{G}''_\eta$ se prolonge en un
$S$-morphisme $v : \underset{\sim}{H} \to \underset{\sim}{G}''$. Comme
$\underset{\sim}{H}$ est connexe et $\underset{\sim}{G}''$ étale, $v$ est nul.
\struck{P} Il en résulte que $u_\eta$ se factorise à travers
$\underset{\sim}{G}'_\eta$, donc toujours d'après Tate, se prolonge en un
$S$-morphisme $u : \underset{\sim}{H} \to \underset{\sim}{G}'$ c.q.f.d.
\marginal{« donc » abusif}
\note{« On sait » est entouré au crayon et relié à la note marginale ;
« un » et « étale » sont écrits à l'encre bleue au-dessus de la ligne, « et »
à l'encre bleue en fin de ligne à la place de « donc » biffé. Le second
« Tate » est souligné deux fois.}

b) Cas général. On sait qu'il existe une extension de traits $S' \to S$,
finie et plate, telle que, si $\eta'$ est le point générique de $S'$, le $S'$
modèle de Néron $\mathcal{G}'$ de $G_\eta$, soit semistable. D'après a), le
morphisme $u_{\eta'} : \underset{\sim}{H}_{\eta'} \to \underset{\sim}{G}_{\eta'}$
se prolonge en un $S'$-morphisme
$u' : \underset{\sim}{H}_{S'} \to \mathcal{G}'$. Rappelons que le $S$-modèle
de Néron de $G_\eta$, c'est à dire $G$, s'obtient à partir de $\mathcal{G}'$
de la façon suivante (cf.\ note sur les modèles de Néron prop.\ 2.5) :
On considère le $S$-groupe $S$-néronien
$\mathcal{G} = \prod_{S'/S} (\mathcal{G}'/S')$. Le groupe $G_\eta$
s'identifie canoniquement à un sous-groupe de $\mathcal{G}_\eta$ ; soit $G_1$
l'adhérence schématique dans $\mathcal{G}$ de $G_\eta$. Alors $G$ est le
« lissifié » de $G_1$.
\marginal{? Cas où $G_\eta$ est une V.A.\ seulement ! OK}
\note{« On sait » est de nouveau entouré au crayon ; la note marginale, au
crayon et encerclée, est en face des premières lignes de b) ; « OK » est
ajouté dessous. Un petit « v » bleu surmonte « semistable ». Le
« $\mathcal{G}'$ » et le « $\prod$ » sont écrits à l'encre bleue dans les
blancs du tapuscrit.}

\page{49}
Ceci étant, il résulte de la définition de $\prod_{S'/S}$ et de l'existence
du prolongement $u'$, que $u_\eta$ se prolonge en un $S$-morphisme
$u_1 : \underset{\sim}{H} \to G_1$. Il reste à voir que $u_1$ se relève en un
$S$-morphisme $u : \underset{\sim}{H} \to G$. Comme $\underset{\sim}{H}$ est
connexe, il est loisible de supposer que $p$ est la caractéristique
résiduelle de $S$. Or on a le :
\note{après « $\underset{\sim}{H} \to G$ », un second $G$ tapé en surcharge ;
les flèches de ce feuillet sont complétées à l'encre bleue.}

\emph{Lemme 3} : Si $f : G \to G_1$ désigne le morphisme canonique, il
existe $n \geqslant 0$, tel que l'élévation à la puissance $p^n$ dans $G_1$ se
factorise en :
\[
  G_1 \xrightarrow{\;g\;} G \xrightarrow{\;f\;} G_1 .
\]
(ceci a été démontré au cours du séminaire de l'an dernier).
\marginal{(qui induit isom.\ sur fibres gén.)}
\note{note au crayon au-dessus de l'énoncé, reliée par une accolade à
« morphisme canonique » ; « isom. » est une lecture probable.}

On a alors le diagramme commutatif suivant, dans le quel la première colonne
est exacte :

\begin{tikzcd}[column sep=small, row sep=small]
0 \arrow[d] & & \\
H_n \arrow[d] & & \\
\underset{\sim}{H} \arrow[dd, "p^n"'] \arrow[rr, "u_1"] & & G_1 \arrow[d, "g"] \arrow[dd, bend left=60, "p^n"] \\
& & G \arrow[d, "f"] \\
\underset{\sim}{H} \arrow[d] \arrow[urr, dashed, "u"] \arrow[rr, "u_1"'] & & G_1 \\
0 & &
\end{tikzcd}

\note{diagramme redessiné. Les flèches verticales de la colonne de gauche sont
à l'encre bleue, la flèche $u$ en pointillé ; l'arc $p^n$ de $G_1$ à $G_1$,
à droite, est au crayon.}

Le morphisme $g\,u_1$ s'annule évidemment sur $H_n$ puisqu'il en est ainsi sur
la fibre générique, donc il se factorise à travers
$\underset{\sim}{H} = \underset{\sim}{H}/H_n$, d'où l'existence du relèvement
$u$.

\emph{Remarques} : La démonstration utilise à deux reprises le théorème
de Tate, qui en toute rigueur n'est démontré que lorsque $\eta$ est de
caractéristique $0$. La première fois, le groupe d'arrivée
$\underset{\sim}{G}''$ est étale et dans ce cas le théorème de Tate n'est pas
très difficile. Par ailleurs, je sais maintenant prouver le théorème de Tate
dans le cas où la partie connexe, sur la fibre générique, dans le groupe
d'arrivée, est de type multiplicatif.

\emph{Corollaire}. Soit $G$ un \struck{$S$} [$\eta$]-groupe
\struck{$S$-néronien,} dont la fibre générique est exten-

\page{50}
sion d'une variété abélienne par un tore. Alors, pour que $G$ \struck{es} ait
« bonne réduction » sur $S$, il faut et il suffit que le groupe
$p$-divisible $\underset{\sim}{G}$ se prolonge en un $S$-groupe
$p$-divisible.
\marginal{de car.\ rés.\ $p > 0$}
\marginal{explicitez ce que ça veut dire}
\note{la première note, au crayon et encerclée en tête du feuillet, est reliée
par un trait à « sur $S$ » ; la seconde, en marge gauche, est reliée par une
accolade à « bonne réduction sur $S$ ». Les guillemets autour de « bonne
réduction » sont ajoutés à la main. Le tapuscrit s'arrête ici ; le reste du
feuillet est vierge.}

\section{[Homomorphismes d'un groupe $p$-divisible dans un groupe algébrique]}
\note{titre de l'éditeur, entre crochets : les feuillets 51 et 52 n'en portent
aucun. Encre bleue, écriture soignée.}

\page{51}
\emph{Lemme}. Soient $K$ un corps, $G$ un groupe algébrique sur $K$,
$M = (M(n))$ un groupe $p$-divisible ($p$ un nb premier
\struck{\uncertain{$\neq$ car $K$}}), $\varphi : M \to G$ un hom de groupes.
Alors il existe un sous-schéma en groupe $H$ de $G$, extension d'une VA par un
tore, tel que $\varphi$ se factorise par $H$.
\struck{b) Si car.\ $K =$}
[On suppose que car $K = p$, ou que $K$ parfait, et que $G$ ne contient pas de
ss-groupe isomorphe à $\mathbf{G}_a$.]
\note{la fin de la parenthèse sur $p$ est biffée ; « Alors » est précédé d'un
trait vertical. L'hypothèse entre crochets est un ajout encadré, relié par
un trait à la fin de l'énoncé, à la place d'un « b) » biffé.}

\emph{Dém}. On a $\varphi = (\varphi(n))$, avec
\[
  \varphi(n) : M(n) \longrightarrow G .
\]
Soit $Z_n = \operatorname{Centr} \varphi(n)$. Les $Z_n$ sont des sous-schémas
en groupe formant une suite décroissante, donc stationnaire, soit $Z$ la
valeur commune. Évidemment $\varphi$ se factorise par \struck{$G$} $Z$, et
même par $\operatorname{Centr}(Z)$. Cela nous ramène au cas où $G$ est
commutatif. Évidemment l'image composée $M \to G \to G/G^{\circ}$ est nulle,
donc on peut supposer $G$ connexe.

\emph{1er Cas} $p = \mathrm{car.}\, K$. On a évidemment, pour tout
$\nu \geqslant 0$, que $\varphi$ se factorise par $p^\nu G$. Or pour $\nu$
grand, on sait \struck{$p^\nu G$} que $p^\nu G$ est une extension d'une VA par
un tore (à dégager en lemme, valable pour tout groupe algébrique commutatif
sur un corps de car.\ $p > 0$ \struck{ou} [ou $K$ \uncertain{parfait}, $G$
n'ayant pas de ss-groupe isom.\ à $\mathbf{G}_a$]).
\note{l'ajout final est écrit au-dessus de la ligne, en plus petit, au-dessus
d'un « ou » biffé ; le mot après $K$ se lit mal.}

\emph{2ème Cas} \struck{$p \neq$ car $K$, les $M(n)$ sont alors
étales.} Alors on sait que $G$ est ext.\ d'une VA.\ par un tore !

\page{52}
\struck{Le groupe \uncertain{engendré} [\uncertain{algébrique}] par les
$\varphi(M(n))$ est donc défini \ill{} un sous-groupe algébrique lisse de $G$.
D\ill{}}
\note{deux lignes encadrées et biffées de traits obliques en tête du
feuillet.}

Supposons maintenant que $G$ est un schéma en groupes Néronien sur un trait
$S$ de corps des fonctions $K$, que $M$ un groupe $p$-divisible, et
$\varphi_K : M_K \to G_K$. Je dis que $\varphi_K$ se prolonge en
$\varphi : M \to G$ (évidemment unique).

Clair si $p \neq$ car.\ résiduelle.

Supposons donc $p =$ car.\ $k$. Alors, ou bien car.\ $K = p$, ou bien car.\
$K = 0$ (donc $K$ parfait) ; en tout cas, $G_K$ n'a pas de ss-groupe isom.\ à
$\mathbf{G}_a$, puisque $G_K$ admet un modèle de Néron. Donc on peut
appliquer le lemme précédent, \struck{Il existe} d'où un $H_K \subset G_K$
par lequel se factorise $\varphi$. Comme $H_K$ est ext.\ d'une VA par un tore,
$H_K$ admet un modèle de Néron $H$, et $H_K \xrightarrow{i_K} G_K$ se
prolonge par $H \xrightarrow{i}$ \uncertain{$K$}. Le th.\ de Raynaud pour $H$
implique que $\psi_K : M_K \to H_K$ se prolonge en $\psi : M \to H$. On gagne !
\note{le but de $i$ est écrit comme un $K$ ; on attend $G$.}

\end{document}
